This section formulates the bi-objective optimization problem for CCHP system capacity planning and describes the solution algorithm.

\subsection{Decision variables}
\label{sec:decision_vars}

The decision vector $\mathbf{x}$ comprises the rated capacities of all system components:
\begin{equation}
    \mathbf{x} = [P_{\text{pv}}, P_{\text{wt}}, P_{\text{gt}}, P_{\text{hp}}, P_{\text{ec}}, P_{\text{ac}}, P_{\text{es}}, P_{\text{hs}}, P_{\text{cs}}]^{\top}
    \label{eq:decision_base}
\end{equation}

When the Carnot battery is enabled, two additional variables are appended:
\begin{equation}
    \mathbf{x}_{\text{CB}} = [P_{\text{pv}}, P_{\text{wt}}, P_{\text{gt}}, P_{\text{hp}}, P_{\text{ec}}, P_{\text{ac}}, P_{\text{es}}, P_{\text{hs}}, P_{\text{cs}}, P_{\text{cb}}, E_{\text{cb}}]^{\top}
    \label{eq:decision_cb}
\end{equation}

All decision variables are continuous and bounded:
\begin{equation}
    0 \leq x_j \leq x_j^{\text{ub}}, \quad j = 1, \ldots, \text{Dim}
    \label{eq:bounds}
\end{equation}
where $x_j^{\text{ub}}$ is the upper bound for the $j$-th variable, determined by site-specific constraints such as available roof area (PV), wind resource (WT), and gas supply capacity (GT). The specific upper bounds for each case study are provided in \Cref{sec:cases}.

\subsection{Objective functions}
\label{sec:objectives}

The bi-objective optimization simultaneously minimizes the annualized total cost and the source-load matching index.

\subsubsection{Objective 1: Annualized total cost}

The first objective aggregates the annualized investment cost, the capacity charge for grid connection, and the annual operational cost:
\begin{equation}
    f_1(\mathbf{x}) = C_{\text{inv}}(\mathbf{x}) + C_{\text{cap}}(\mathbf{x}) + C_{\text{op}}(\mathbf{x})
    \label{eq:obj1}
\end{equation}

The annualized investment cost is computed as:
\begin{equation}
    C_{\text{inv}}(\mathbf{x}) = \sum_{j=1}^{\text{Dim}} \alpha_j \cdot x_j
    \label{eq:invest}
\end{equation}
where $\alpha_j$ is the annualized investment coefficient for the $j$-th component, incorporating the capital cost, the capital recovery factor (based on discount rate and equipment lifetime), and the annual maintenance cost.

The capacity charge reflects the demand charge imposed by the utility based on the annual peak grid power import:
\begin{equation}
    C_{\text{cap}}(\mathbf{x}) = c_{\text{cap}} \cdot \max_{t \in [1,T]} P_{\text{grid}}(t)
    \label{eq:cap_charge}
\end{equation}
where $c_{\text{cap}}$ is the capacity charge rate (\si{\text{\euro}/(kW \cdot a)} or \si{\text{\textyen}/(kW \cdot a)}) and $P_{\text{grid}}(t)$ is the grid import power at hour $t$, obtained from the operational dispatch solution.

The annual operational cost $C_{\text{op}}(\mathbf{x})$ is computed via the typical-day decomposition described in \Cref{sec:typical_day}, summing the weighted daily costs of electricity purchase and gas consumption.

\subsubsection{Objective 2: Source-load matching index}

The second objective is the energy-quality-weighted Euclidean distance matching index defined in \Cref{eq:mi_eqd}:
\begin{equation}
    f_2(\mathbf{x}) = \text{MI}_{\text{EQD}}(\mathbf{x})
    \label{eq:obj2}
\end{equation}

For comparative experiments, $f_2$ is replaced by the corresponding alternative index ($\text{MI}_{\text{Std}}$, $\text{MI}_{\text{Pearson}}$, or $\text{MI}_{\text{SSR}}$) while keeping all other aspects of the framework identical. In the single-objective economic case, $f_2$ is omitted entirely.

\subsection{Optimization algorithm}
\label{sec:algorithm}

The bi-objective optimization problem is solved using the Non-dominated Sorting Genetic Algorithm II (NSGA-II) \cite{deb2002fast}, implemented via the geatpy evolutionary computation library. NSGA-II is chosen for its proven effectiveness in generating well-distributed Pareto fronts for engineering design problems with two to three objectives.

The algorithmic workflow is as follows:

\begin{enumerate}
    \item \textbf{Initialization}: A population of $N_{\text{ind}}$ individuals is generated by Latin hypercube sampling within the feasible bounds. For methods that benefit from warm-starting, the final population of a previously solved related problem can be inherited as the initial population.
    \item \textbf{Fitness evaluation}: For each individual $\mathbf{x}$ in the population, the operational dispatch model (\Cref{sec:dispatch}) is solved for all 14 typical days. The two objective values $f_1(\mathbf{x})$ and $f_2(\mathbf{x})$ are computed from the dispatch results. This step is parallelized across individuals using a process pool to exploit multi-core processors.
    \item \textbf{Selection and reproduction}: NSGA-II performs non-dominated sorting and crowding distance assignment to rank individuals. Tournament selection, simulated binary crossover (SBX), and polynomial mutation generate offspring.
    \item \textbf{Termination}: The algorithm runs for $G_{\text{max}}$ generations. The final non-dominated set constitutes the Pareto front, representing the trade-off between cost and matching performance.
\end{enumerate}

For the single-objective economic optimization, a standard genetic algorithm with differential evolution (DE) mutation is used instead, as the problem reduces to a scalar minimization.

\subsection{Population inheritance mechanism}
\label{sec:inheritance}

To improve convergence efficiency when multiple matching methods are compared on the same case, a population inheritance mechanism is employed. The final population from the standard deviation method (Std) is used as the initial population for the Euclidean distance method (EQD). Since both methods share the same decision space and the first objective (cost), the inherited population provides a warm start that accelerates convergence toward the Pareto front of the new matching index. This mechanism does not bias the final result, as NSGA-II's selection pressure ensures that the population evolves toward the true Pareto front of the target problem regardless of initialization.

\subsection{Computational implementation}
\label{sec:implementation}

The entire framework is implemented in Python. The operational dispatch model is built using oemof-solph and solved by Gurobi (with GLPK as fallback). The evolutionary algorithm uses geatpy with a process-based parallel pool for fitness evaluation. Each fitness evaluation involves solving 14 linear programs (one per typical day), taking approximately 0.5--2 seconds per individual depending on system complexity. A typical optimization run with $N_{\text{ind}} = 50$ individuals and $G_{\text{max}} = 100$ generations requires approximately 2--4 hours on a workstation with 8 CPU cores.
